\section{$G=\Gm$}\srcpage{33}

本节把一切都对 $G=\Gm$ 的情形显式化。同样，我们假定 $X$ 的亏格 $g$ 满足 $g\ge 1$。为简单起见，我们还将固定一个基点 $x\in X$。此时 $\Pic=\Bun_{\Gm}$ 由
\[ \Pic=\Z\times J\times */\Gm \]
给出，其中 $J=\Pic^{0}$ 是 $X$ 的 Jacobian，一个主极化阿贝尔簇。因子 $*/\Gm$ 来自考虑 Picard 叠——若像讨论 Picard 概形时惯常所做的那样，在 $x\in X$ 处把线丛平凡化，我们得到的就会是 $\Z\times J$。具有幂零奇异支撑的条件，就简单地变成具有零奇异支撑的条件，而后者等价于该层是局部常值的。

因此，在 Betti 情形，我们得到
\[ D_{\Nilp}(\Pic)=\prod_{\Z}\Ind\Lisse(J)\otimes D(*/\C^{\times})=\prod_{\Z}D(*/H_{1}(X))\otimes D(*/\C^{\times}) \]
因为 $\pi_{1}(J)=H_{1}(X)=H_{1}(X,\Z)\cong\Z^{2g}$ 是 $X$ 的一阶同调群（$\Z$ 系数）。

\srcpage{34}在对偶一侧，我们计算 $\LocSys^{\Betti}_{\Gm}$。这是把任意 $k$-代数 $A$ 送到（作为群的）映射 $\pi_{1}\to A^{\times}$ 的叠，模去 $A^{\times}$ 经由共轭的平凡作用。正如我们在第 4 节对一般群所讨论的，
\[ \Hom(\pi_{1},\Gm)=\Gm^{2g}\times^{L}_{\Gm}*, \]
其中纤维积在导出意义下取。实际上，由于 $\Gm$ 交换，映射 $\Gm^{2g}\to\Gm$ 恒等于 $1$，于是
\[ \Hom(\pi_{1},\Gm)=\Gm^{2g}\times(*\times^{L}_{\Gm}*), \]
其中 $*\times^{L}_{\Gm}*=\Spec(\e[\epsilon])$ 是相应于 $\e[\epsilon]$ 的仿射导出概形，此处 $\epsilon$ 的次数为 $1$（因而平方为零，所以 $\e[\epsilon]=\e\oplus\e[1]$ 是平方零扩张）。总之，
\[ \LocSys^{\Betti}_{\Gm}=\Gm^{2g}\times(*\times^{L}_{\Gm}*)\times*/\Gm. \]

我们现在可以手算看出存在等价
\[ D_{\Nilp}(\Pic)\cong D_{\qc}(\LocSys^{\Betti}_{\Gm}). \]
事实上，这是三个等价之张量积：

\begin{enumerate}
\item 第一，
\[ \prod_{\Z}D(*)=D_{\qc}(*/\Gm): \]
这来自按权 $n\in\Z$ 对 $\Gm$-表示的分解。

\item 第二，
\[ D(*/H_{1}(X))\cong D_{\qc}(\Gm^{2g}): \]
事实上，$H_{1}(X)\cong\Z^{2g}$ 的一个表示是一个带有 $2g$ 个交换自同态的向量空间，它同时也是 $\Gm^{2g}$ 上的拟凝聚层。

\item 第三，
\[ D(*/\C^{\times})\cong D_{\qc}(\Spec(\e[\epsilon])). \]
事实上，$D(*/\C^{\times})$ 由单位对象沿光滑映射 $*\to*/\C^{\times}$ 的 $!$-推前紧生成，而该对象的自同态可算出是 $\C^{\times}$ 的同调，即 $\e[\epsilon]$。
\end{enumerate}

于是，在模式匹配的层面上，就得到所要的等价。我们注意到，对第三部分，我们需要 $\Spec(\e[\epsilon])$ 这一项，它要求人在导出代数几何的框架中工作。稍加仔细地做，可以证明谱作用把 $D_{\Nilp}(\Pic)$ 变成 $D_{\qc}(\LocSys^{\Betti}_{\Gm})$ 上的自由秩 $1$ 模。

\begin{remark}
上述所有等价都落入「Cartier 对偶」的一般模式，而在交换情形，Langlands 对偶确实可以看作 Cartier 对偶的一个应用。粗略地说，对任意类型的交换群概形或叠 $G$，都应有由（导出的）$\Hom(G,\Gm)$ 给出的 Cartier 对偶 $G^{*}$，此时应有等价 $D_{\qc}(*/G)\cong D_{\qc}(G^{*})$ 与 $D_{\qc}(*/G^{*})\cong D_{\qc}(G)$。例如，这出现在有限平坦群概形的经典情形中，\srcpage{35}但它对 $G=\Z$ 与 $G^{*}=\Gm$ 也成立。在阿贝尔簇 $A$ 的情形，对偶实际上是对偶阿贝尔簇 $A^{\vee}[-1]$，只是作了上同调次数 $1$ 的平移；此时该等价就变成 Fourier--Mukai 等价 $D_{\qc}(A)\cong D_{\qc}(A^{\vee})$。在 Gestalten 的框架中存在一般的 Cartier 对偶理论，参见 Scholze（2026）。
\end{remark}

在 étale 情形，故事以非常平行的方式展开。让我们更详细地讨论 de Rham 情形。在一侧，我们得到 $\Pic$ 上的 D-模范畴。同样，分解
\[ \Pic=\Z\times J\times */\Gm \]
把 $D(\Pic)$ 分解为三个因子 $D(\Z)$、$D(J)$ 与 $D(*/\Gm)$ 的张量积。因子 $D(\Z)$ 与 $D(*/\Gm)$ 等价于它们在 Betti 情形中的版本，并在 Langlands 对偶一侧以平行方式处理。然而 $D(J)=D_{\qc}(J_{\dR})$ 与它的 Betti 对应物颇不相同。

在 Langlands 对偶一侧，$\LocSys^{\dR}_{\Gm}$ 映射到 $\Pic$，而纤维由 $\Gamma(X,\OmX)$ 下的一个拟挠子给出。实际上，恰好命中 $\Pic$ 的次数 $0$ 部分 $J\times*/\Gm$。把一切展开，并同时记下 $H^{1}(X,\OmX)=k$，可以看出
\[ \LocSys^{\dR}_{\Gm}=EJ\times*/\Gm\times\Spec(k[\epsilon]), \]
其中
\[ 0\to V J\to EJ\to J\to 0 \]
是 Jacobian $J$ 的通用向量扩张，其核是对应于向量空间 $H^{0}(X,\OmX)$ 的向量群 $V J$。

现在结果表明，$J_{\dR}=J/(0\subset J)^{\wedge}$ 的 Cartier 对偶恰是 $EJ[-1]$——由于 $J$ 主极化，$J$ 本身对偶到 $J[-1]$，而形式交换群 $(0\subset J)^{\wedge}$ 对偶到 $V J$（因为一般地形式加法群是加法群的 Cartier 对偶）。因此 $D_{\qc}(J_{\dR})\cong D_{\qc}(EJ)$——这一 Cartier 对偶的实例归功于 Laumon（1996）。

注意 $\LocSys^{\Betti}_{\Gm}$ 中的因子 $\Gm^{2g}$ 被 $\LocSys^{\dR}_{\Gm}$ 中的因子 $EJ$ 替换。正如 Riemann--Hilbert 对应所预言的，它们有相同的 $\C$-点，甚至在 $\C$-点处有相同的形式完备化，但作为代数簇却非常不同。最重要的是，$\Gm^{2g}$ 是仿射的，因而有许多整体函数，而 $EJ$ 的整体截面只有 $k$。de Rham 情形中整体截面的这种稀少，实际上在几何 Langlands 定理证明的最后一步起了关键作用。
